\chapter{هم‌ارزی محرمانگی تفاضلی موضعی و انقباض $f$-واگرایی‌ها}
\label{ch:ldp-contraction}

\section{مقدمه}
\label{sec:ldp-contraction:intro}

در فصل پیشین (\ref{ch:minimax-theory})، چارچوب نظریه مینی‌مکس بر اساس کارهای بنیادین دوچی و همکاران \cite{duchi2018} را بررسی کردیم. مشاهده شد که اعمال قید محرمانگی تفاضلی موضعی (\lr{LDP}) به صورت ذاتی موجب کاهش دقت آماری می‌شود و نرخ همگرایی خطا را از مرتبه $\mathcal{O}\left( \frac{1}{n} \right)$ در مدل متمرکز به مرتبه $\mathcal{O}\left( \frac{1}{n\al^2} \right)$ تنزل می‌دهد. رویکرد دوچی برای اثبات این کران‌های پایین، بر پایه تحدید فواصل اطلاعاتی میان خروجی‌های مکانیزم \mech\ استوار بود. در آن چارچوب، با استفاده از نامساوی‌های خاصی، واگرایی \lr{KL} و اطلاعات متقابل خروجی‌ها بر حسب فاصله‌ی \lr{TV} توزیع‌های ورودی و بودجه محرمانگی $\al$ کران‌دار می‌شدند و سپس این کران‌ها به عنوان ورودی به نامساوی‌های فانو و لوکم تغذیه می‌شدند. اگرچه این روش کران‌های مینی‌مکس بهینه‌ای از منظر مرتبه مجانبی ارائه می‌دهد، اما ماهیت آن بر پایه نامساوی‌های موردی است و ساختار هندسی دقیقِ تبدیل اطلاعات در کانال‌های خصوصی را به طور کامل مشخص نمی‌کند. 

در این فصل، یک تغییر رویه‌ی اساسی در تحلیل \lr{LDP} ارائه می‌دهیم. با الهام از پژوهش آسوده و همکاران \cite{Asoodeh2021}، نشان می‌دهیم که قید محرمانگی تفاضلی موضعی صرفاً یک محدودیت الگوریتمی نیست، بلکه به طور دقیق و ریاضیاتی معادل با ویژگی «انقباض» در یک رده خاص از فواصل توزیعی به نام $E_\gamma$-واگرایی است. این هم‌ارزی نظری، پیوندی عمیق و بنیادین میان نظریه حریم خصوصی و نامساوی‌های قوی پردازش داده برقرار می‌سازد. در این نگاه جدید، هر مکانیزم محرمانگی موضعی \mech\ به عنوان یک کانال نویزی تفسیر می‌شود که توانایی تفکیک‌پذیری آماری بین توزیع‌ها را با یک ضریب مشخص منقبض می‌کند.

تفاوت بنیادین این رویکرد با کارهای پیشین در جامعیت، دقت و قابلیت تعمیم آن است. در چارچوب دوچی، برای هر کلاس از مسائل (مانند تخمین میانگین یا آزمون فرض) نیازمند اثبات نامساوی‌های مجزایی برای اتصال قید حریم خصوصی به ابزارهای نظریه تصمیم بودیم. اما با استفاده از مفهوم ضریب انقباض، می‌توانیم تأثیر مکانیزم \mech\ را به صورت یک ضریب ضرب‌شونده‌ی دقیق بر روی هر $f$-واگرایی دل‌خواه اعمال کنیم. به عبارت دیگر، با اثبات هم‌ارزی \lr{LDP} با انقباض $E_\gamma$-واگرایی، می‌توانیم از طریق نمایش انتگرالی ساسون (رابطه \ref{eq:integral_rep_sason} از فصل \ref{ch:preliminaries})، ضرایب انقباض بهینه را برای سایر اعضای خانواده $f$-واگرایی‌ها (مانند \lr{KL, $\chi^2, H^2$}) استخراج نماییم.

علاوه بر این، اگر رده‌ی تمامی مکانیزم‌های تصادفی که قید محرمانگی تفاضلی موضعی تقریبی با پارامترهای $\al$ و $\del$ را برآورده می‌کنند با نماد $\mathcal{Q}_{\al,\del}$ نشان دهیم، کارهای دوچی عمدتاً محدود به تحلیل مکانیزم‌های متعلق به رده‌ی محرمانگی محض، یعنی $\mech \in \mathcal{Q}_{\al,0}$، است. در مقابل، چارچوب نوین مبتنی بر ضریب انقباض، به صورت تحلیلی و یک‌پارچه به کلاس کلی‌تر $\mathcal{Q}_{\al,\del}$ نیز بسط می‌یابد و ابزاری قدرتمند برای تحلیل تمامی مکانیزم‌های عضو این رده فراهم می‌کند. از منظر دقت کران‌های به‌دست‌آمده نیز، مقایسه‌های انجام‌شده در مقاله آسوده نشان می‌دهد که کران‌های استخراج‌شده مبتنی بر ضریب انقباض برای مکانیزم‌های $\mech \in \mathcal{Q}_{\al,\del}$ در برخی موارد دقیق‌تر از کران‌های دوچی بوده و بهبودی در حد یک ضریب ثابت ایجاد می‌کنند، در حالی که در بعضی موارد دیگر ممکن است اندکی بدتر عمل کنند. با این وجود، این دستاورد یک‌پارچه، نه تنها اثبات کران‌های پایین مینی‌مکس را به شدت نظام‌مند می‌سازد، بلکه بینش هندسی عمیق‌تری از ماهیت تبادل میان محرمانگی و سودمندی به دست می‌دهد و نشان می‌دهد که افت کارایی ناشی از محرمانگی موضعی، یک سد اطلاعاتی-هندسی غیرقابل عبور است.
\section{ضریب انقباض و تحلیل هندسی}
\label{sec:contraction-geometry}

برای کمّی‌سازی میزان کاهش فاصله توزیع‌ها پس از عبور از یک کانال نویزی (مکانیزم \mech)، از مفهوم ضریب انقباض در هندسه اطلاعات استفاده می‌شود. این ضریب به ما نشان می‌دهد که یک عملگر انتقال مارکوفی تا چه حد فضای توزیع‌های احتمال را منقبض می‌کند و تفکیک‌پذیری آماری آن‌ها را تقلیل می‌دهد.


\subsection{کران‌های عمومی انقباض برای $f$-واگرایی‌ها (لم ۱ و ۲)}

با اثبات هم‌ارزی در قضیه \ref{thm:ldp-equivalence}، مقاله آسوده نشان می‌دهد که شرط محرمانگی موضعی به شکل یکتا می‌تواند ضریب انقباض را برای \textbf{تمامی} اعضای خانواده $f$-واگرایی‌ها کران‌دار کند. این دستاورد کلیدی در قالب لم ۱ ارائه شده است.

\begin{لم}[کران عمومی ضریب انقباض]
\label{lem:general-contraction}
فرض کنید $\mech \in \mathcal{Q}_{\al,\del}$ یک مکانیزم محرمانگی تفاضلی موضعی تقریبی باشد. اگر تابع $\vph(\al, \del)$ را به صورت زیر تعریف کنیم:
\begin{equation}
\label{eq:phi-def}
\vph(\al, \del) := 1 - (1 - \del)e^{-\al}
\end{equation}
در این صورت، برای هر $f$-واگرایی دل‌خواه، ضریب انقباض مکانیزم $\mech$ توسط این تابع کران‌دار می‌شود؛ یعنی $\eta_f(\mech) \lee \vph(\al, \del)$. به بیان دیگر، برای هر دو توزیع ورودی دل‌خواه $P$ و $Q$ روی فضای \Xset، نامساوی زیر همواره برقرار است:
\begin{equation}
\label{eq:lem1-asoodeh}
D_f\left( \mech(P) \| \mech(Q) \right) \lee \vph(\al, \del) D_f(P \| Q)
\end{equation}
\end{لم}

\begin{اثبات}
بر اساس یک نتیجه کلاسیک در نظریه اطلاعات (معروف به قضیه آلسوِد-گاکس\LTRfootnote{Ahlswede-Gács Theorem})، می‌دانیم که ضریب انقباض برای هر $f$-واگرایی دل‌خواه هرگز نمی‌تواند از ضریب انقباض متغیرات کل (که به ضریب دوبروشین\LTRfootnote{Dobrushin Contraction Coefficient} معروف است) فراتر رود. به بیان ریاضی:
\begin{equation}
\label{eq:lem1-proof-1}
\eta_f(\mech) \lee \eta_{TV}(\mech) = \sup_{x, x' \in \Xset} d_{TV}\left( \mech(x), \mech(x') \right)
\end{equation}
بنابراین، برای اثبات لم کافیست بیشینه فاصله $TV$ بین خروجی‌های هر دو نقطه ورودی $x$ و $x'$ را کران‌دار کنیم. فرض کنید $p = \mech(\cdot|x)$ و $q = \mech(\cdot|x')$. می‌دانیم که فاصله متغیرات کل را می‌توان بر حسب پیشامدی که بیشترین اختلاف را ایجاد می‌کند نوشت: $d_{TV}(p, q) = p(S) - q(S)$ که در آن $S = \left\{ z \in \Zset : p(z) > q(z) \right\}$.

از آنجا که مکانیزم $\mech$ در شرط \LDP[(\al,\del)] صدق می‌کند، برای پیشامد $S$ داریم:
\begin{equation}
\label{eq:lem1-proof-2}
p(S) \lee e^\al q(S) + \del \implies q(S) \gee e^{-\al} \left( p(S) - \del \right)
\end{equation}
با جایگذاری نامساوی \eqref{eq:lem1-proof-2} در رابطه فاصله $TV$، به دست می‌آوریم:
\begin{align}
d_{TV}(p, q) &= p(S) - q(S) \nonumber \\
&\lee p(S) - e^{-\al} \left( p(S) - \del \right) \nonumber \\
&= p(S) \left( 1 - e^{-\al} \right) + \del e^{-\al} \label{eq:lem1-proof-3}
\end{align}
از آنجا که $p(S)$ یک مقدار احتمال است و همواره $p(S) \lee 1$ برقرار است، و همچنین با توجه به اینکه عبارت $\left( 1 - e^{-\al} \right)$ مقداری نامنفی است، بیشترین مقدار رابطه \eqref{eq:lem1-proof-3} زمانی رخ می‌دهد که $p(S) = 1$ باشد. در نتیجه:
\begin{equation}
\label{eq:lem1-proof-4}
d_{TV}(p, q) \lee 1 \left( 1 - e^{-\al} \right) + \del e^{-\al} = 1 - e^{-\al} + \del e^{-\al} = 1 - e^{-\al} (1 - \del)
\end{equation}
سمت راست رابطه \eqref{eq:lem1-proof-4} دقیقاً برابر با تعریف $\vph(\al, \del)$ است. با ترکیب این نتیجه با رابطه \eqref{eq:lem1-proof-1}، اثبات لم به پایان می‌رسد.
\end{اثبات}

این لم بسیار بنیادین است، زیرا نشان می‌دهد که مستقل از انتخاب نوع فاصله اطلاعاتی (مانند واگرایی \lr{KL}، $\chi^2$، $H^2$ یا $TV$)، میزان حفظ اطلاعات در خروجی یک مکانیزم \LDP[(\al,\del)] هرگز نمی‌تواند از ضریب $\vph(\al, \del)$ فراتر رود. به عنوان نمونه، در حالت محرمانگی محض ($\del = 0$)، این ضریب استهلاک اطلاعات دقیقاً برابر با $1 - e^{-\al}$ خواهد بود.

در تحلیل مینی‌مکس و نظریه اطلاعات (که هدف نهایی ما برای اثبات کران‌های پایین است)، ما غالباً با مجموعه‌ای از $n$ کاربر مواجهیم که هر یک دارای داده‌ی مستقل $X_i$ هستند و هر کدام داده‌ی خود را به صورت مستقل از یک کانال خصوصی عبور می‌دهند. لم دوم مقاله، این کران انقباض را برای چنین ترکیب مستقلی (فضای ضربی) تعمیم می‌دهد.

\begin{لم}[انقباض برای ترکیبات مکانیزم]
\label{lem:tensor-contraction}
فرض کنید $\mech \in \mathcal{Q}_{\al,\del}$ و عملگر $\mech^{\otimes n}$ نشان‌دهنده اعمال مستقل این مکانیزم بر روی بردار $n$تایی داده‌های ورودی باشد. با تعریف تابع $\vph_n(\al, \del)$ به صورت:
\begin{equation}
\label{eq:phi-n-def}
\vph_n(\al, \del) := 1 - e^{-n\al}(1 - \del)^n
\end{equation}
برای هر $n \gee 1$ و هر $f$-واگرایی، ضریب انقباض این مکانیزم ترکیبی در نامساوی زیر صدق می‌کند:
\begin{equation}
\label{eq:lem2-asoodeh}
\eta_f\left( \mech^{\otimes n} \right) \lee \vph_n(\al, \del)
\end{equation}
\end{لم}

\begin{اثبات}
همانند اثبات لم پیشین، ضریب انقباض کل $\eta_f\left( \mech^{\otimes n} \right)$ توسط ضریب دوبروشین (بیشینه فاصله $TV$ بین خروجی‌های هر دو بردار ورودی $X^n$ و $Y^n$) کران‌دار می‌شود.
برای اندازه‌های ضربی\LTRfootnote{Product Measures} (توزیع‌های مستقل)، یک ویژگی شناخته‌شده از فاصله متغیرات کل این است که مکمل آن به صورت ضربی عمل می‌کند:
\begin{equation}
\label{eq:lem2-proof-1}
1 - d_{TV}\left( \mech^{\otimes n}(X^n), \mech^{\otimes n}(Y^n) \right) \gee \prod_{i=1}^{n} \left( 1 - d_{TV}\left( \mech(X_i), \mech(Y_i) \right) \right)
\end{equation}
(دلیل این امر آن است که فاصله $TV$ معادل است با $1 - \int \min(dP, dQ)$، و انتگرال مینیمم برای توزیع‌های مستقل به انتگرالِ ضربِ مینیمم‌ها تقلیل می‌یابد).
از لم \ref{lem:general-contraction} می‌دانیم که برای هر مؤلفه $i$، فاصله متغیرات کل همواره توسط $\vph(\al, \del)$ کران‌دار است:
\begin{equation*}
TV\left( \mech(X_i), \mech(Y_i) \right) \lee \vph(\al, \del) \implies 1 - TV\left( \mech(X_i), \mech(Y_i) \right) \gee 1 - \vph(\al, \del)
\end{equation*}
با جایگذاری این کران در رابطه \eqref{eq:lem2-proof-1} نتیجه می‌شود:
\begin{equation}
\label{eq:lem2-proof-2}
1 - d_{TV}\left( \mech^{\otimes n}(X^n), \mech^{\otimes n}(Y^n) \right) \gee \left( 1 - \vph(\al, \del) \right)^n
\end{equation}
با مرتب‌سازی رابطه فوق بر حسب $d_{TV}$ داریم:
\begin{equation}
\label{eq:lem2-proof-3}
d_{TV}\left( \mech^{\otimes n}(X^n), \mech^{\otimes n}(Y^n) \right) \lee 1 - \left( 1 - \vph(\al, \del) \right)^n
\end{equation}
حال کافیست تعریف $\vph(\al, \del) = 1 - e^{-\al}(1 - \del)$ را در داخل پرانتز قرار دهیم:
\begin{align}
1 - \left( 1 - \left( 1 - e^{-\al}(1 - \del) \right) \right)^n &= 1 - \left( e^{-\al}(1 - \del) \right)^n \nonumber \\
&= 1 - e^{-n\al}(1 - \del)^n \label{eq:lem2-proof-4}
\end{align}
عبارت به دست آمده دقیقاً همان تعریف $\vph_n(\al, \del)$ است که اثبات لم را کامل می‌کند.
\end{اثبات}

\textbf{تفسیر هندسی لم‌ها:} این دو لم، هسته‌ی مرکزی اتصال حریم خصوصی به کران‌های پایین مینی‌مکس هستند. برخلاف روش‌های موردی پیشین که نیازمند اثبات‌های مجزا و بسط تیلور مختص به واگرایی \lr{KL} بودند، لم‌های ۱ و ۲ به صورت تحلیلی اثبات می‌کنند که اعمال هر مکانیزم $\mathcal{Q}_{\al,\del}$، فضای تمامی توزیع‌های آماری را با یک ضریب مشخص فشرده می‌سازد. این محدودیت فشردگی مستقیماً به عنوان ورودی به نامساوی‌های فانو و لوکم (که در فصل \ref{ch:preliminaries} بررسی شدند) تغذیه می‌شود و با اعمال آن می‌توانیم افت دقت آماری را با رویکردی کاملاً یکپارچه و جهان‌شمول استخراج نماییم.